% Môn: Toán
% Lớp: 10
% Chương: Chương I. Mệnh đề và tập hợp
% Bài:  Mệnh đề · Mệnh đề với mọi, tồn tại
% Số câu: 0
% App đọc file này trực tiếp — sửa rồi Commit trên GitHub
% Ghi nguồn từng câu: \nguon{SGK} hoặc \nguon{2 SBT VL 12 KNTT} trong


%%%=============EX_1=============%%%

% Mức: NB
% ID: T10C1B1-07-TN
% ID: T10C1B1-07-TN
% Mức: NB
% ID: T10C1B1-06-TN
% ID: T10C1B1-06-TN
% Mức: NB
% ID: T10C1B1-06-TN

% Mức: NB
% ID: T10C1B1-08-TN
% ID: T10C1B1-08-TN
% Mức: NB
% ID: T10C1B1-07-TN
% ID: T10C1B1-07-TN
% Mức: NB
% ID: T10C1B1-07-TN

% Mức: NB
% ID: T10C1B1-09-TN
% ID: T10C1B1-09-TN
% Mức: NB
% ID: T10C1B1-08-TN
% ID: T10C1B1-08-TN
% Mức: NB
% ID: T10C1B1-08-TN

% ID: 0D1N1-1
% Mức: NB
% ID: T10C1B1-10-TN
% ID: T10C1B1-10-TN
% Mức: NB
% ID: T10C1B1-09-TN
% ID: T10C1B1-09-TN
% Mức: NB
% ID: T10C1B1-09-TN

% Mức: NB
% ID: T10C1B1-11-TN
% ID: T10C1B1-11-TN
% Mức: NB
% ID: T10C1B1-10-TN
% ID: T10C1B1-10-TN
% Mức: NB
% ID: T10C1B1-10-TN

% Mức: NB
% ID: T10C1B1-12-TN
% ID: T10C1B1-12-TN
% Mức: NB
% ID: T10C1B1-11-TN
% ID: T10C1B1-11-TN
% Mức: NB
% ID: T10C1B1-11-TN

% Mức: TH
% ID: T10C1B1-13-TN
% ID: T10C1B1-13-TN
% Mức: TH
% ID: T10C1B1-12-TN
% ID: T10C1B1-12-TN
% Mức: TH
% ID: T10C1B1-12-TN

% Mức: NB
% ID: T10C1B1-14-TN
% ID: T10C1B1-14-TN
% Mức: NB
% ID: T10C1B1-13-TN
% ID: T10C1B1-13-TN
% Mức: NB
% ID: T10C1B1-13-TN

% Mức: NB
% ID: T10C1B1-15-TN
% ID: T10C1B1-15-TN
% Mức: NB
% ID: T10C1B1-14-TN
% ID: T10C1B1-14-TN
% Mức: NB
% ID: T10C1B1-14-TN

% Mức: TH
% ID: T10C1B1-16-TN
% ID: T10C1B1-16-TN
% Mức: TH
% ID: T10C1B1-15-TN
% ID: T10C1B1-15-TN
% Mức: TH
% ID: T10C1B1-15-TN

% Mức: TH
% ID: T10C1B1-17-TN
% ID: T10C1B1-17-TN
% Mức: TH
% ID: T10C1B1-16-TN
% ID: T10C1B1-16-TN
% Mức: TH
% ID: T10C1B1-16-TN

% Mức: NB
% ID: T10C1B1-18-TN
% ID: T10C1B1-18-TN
% Mức: NB
% ID: T10C1B1-17-TN
% ID: T10C1B1-17-TN
% Mức: NB
% ID: T10C1B1-17-TN

% Mức: NB
% ID: T10C1B1-19-TN
% ID: T10C1B1-19-TN
% Mức: NB
% ID: T10C1B1-18-TN
% ID: T10C1B1-18-TN
% Mức: NB
% ID: T10C1B1-18-TN

% Mức: NB
% ID: T10C1B1-20-TN
% ID: T10C1B1-20-TN
% Mức: NB
% ID: T10C1B1-19-TN
% ID: T10C1B1-19-TN
% Mức: NB
% ID: T10C1B1-19-TN

% Mức: NB
% ID: T10C1B1-21-TN
% ID: T10C1B1-21-TN
% Mức: NB
% ID: T10C1B1-20-TN
% ID: T10C1B1-20-TN
% Mức: NB
% ID: T10C1B1-20-TN

% Mức: NB
% ID: T10C1B1-22-TN
% ID: T10C1B1-22-TN
% Mức: NB
% ID: T10C1B1-21-TN
% ID: T10C1B1-21-TN
% Mức: NB
% ID: T10C1B1-21-TN

% Mức: NB
% ID: T10C1B1-23-TN
% ID: T10C1B1-23-TN
% Mức: NB
% ID: T10C1B1-22-TN
% ID: T10C1B1-22-TN
% Mức: NB
% ID: T10C1B1-22-TN

% Mức: NB
% ID: T10C1B1-24-TN
% ID: T10C1B1-24-TN
% Mức: NB
% ID: T10C1B1-23-TN
% ID: T10C1B1-23-TN
% Mức: NB
% ID: T10C1B1-23-TN

% Mức: NB
% ID: T10C1B1-25-TN
% ID: T10C1B1-25-TN
% Mức: NB
% ID: T10C1B1-24-TN
% ID: T10C1B1-24-TN
% Mức: NB
% ID: T10C1B1-24-TN

% Mức: NB
% ID: T10C1B1-26-TN
% ID: T10C1B1-26-TN
% Mức: NB
% ID: T10C1B1-25-TN
% ID: T10C1B1-25-TN
% Mức: NB
% ID: T10C1B1-25-TN

% Mức: NB
% ID: T10C1B1-27-TN
% ID: T10C1B1-27-TN
% Mức: NB
% ID: T10C1B1-26-TN
% ID: T10C1B1-26-TN
% Mức: NB
% ID: T10C1B1-26-TN

% Mức: NB
% ID: T10C1B1-28-TN
% ID: T10C1B1-28-TN
% Mức: NB
% ID: T10C1B1-27-TN
% ID: T10C1B1-27-TN
% Mức: NB
% ID: T10C1B1-27-TN

% Mức: NB
% ID: T10C1B1-29-TN
% ID: T10C1B1-29-TN
% Mức: NB
% ID: T10C1B1-28-TN
% ID: T10C1B1-28-TN
% Mức: NB
% ID: T10C1B1-28-TN

% Mức: NB
% ID: T10C1B1-30-TN
% ID: T10C1B1-30-TN
% Mức: NB
% ID: T10C1B1-29-TN
% ID: T10C1B1-29-TN
% Mức: NB
% ID: T10C1B1-29-TN

% Mức: NB
% ID: T10C1B1-31-TN
% ID: T10C1B1-31-TN
% Mức: NB
% ID: T10C1B1-30-TN
% ID: T10C1B1-30-TN
% Mức: NB
% ID: T10C1B1-30-TN

% Mức: NB
% ID: T10C1B1-32-TN
% ID: T10C1B1-32-TN
% Mức: NB
% ID: T10C1B1-31-TN
% ID: T10C1B1-31-TN
% Mức: NB
% ID: T10C1B1-31-TN

% Mức: NB
% ID: T10C1B1-33-TN
% ID: T10C1B1-33-TN
% Mức: NB
% ID: T10C1B1-32-TN
% ID: T10C1B1-32-TN
% Mức: NB
% ID: T10C1B1-32-TN

% Mức: NB
% ID: T10C1B1-34-TN
% ID: T10C1B1-34-TN
% Mức: NB
% ID: T10C1B1-33-TN
% ID: T10C1B1-33-TN
% Mức: NB
% ID: T10C1B1-33-TN

% Mức: NB
% ID: T10C1B1-35-TN
% ID: T10C1B1-35-TN
% Mức: NB
% ID: T10C1B1-34-TN
% ID: T10C1B1-34-TN
% Mức: NB
% ID: T10C1B1-34-TN

% Mức: NB
% ID: T10C1B1-36-TN
% ID: T10C1B1-36-TN
% Mức: NB
% ID: T10C1B1-35-TN
% ID: T10C1B1-35-TN
% Mức: NB
% ID: T10C1B1-35-TN

% Mức: NB
% ID: T10C1B1-37-TN
% ID: T10C1B1-37-TN
% Mức: NB
% ID: T10C1B1-36-TN
% ID: T10C1B1-36-TN
% Mức: NB
% ID: T10C1B1-36-TN

% ID: 0D1N1-1
% Mức: NB
% ID: T10C1B1-38-TN
% ID: T10C1B1-38-TN
% Mức: NB
% ID: T10C1B1-37-TN
% ID: T10C1B1-37-TN
% Mức: NB
% ID: T10C1B1-37-TN

% Mức: NB
% ID: T10C1B1-39-TN
% ID: T10C1B1-39-TN
% Mức: NB
% ID: T10C1B1-38-TN
% ID: T10C1B1-38-TN
% Mức: NB
% ID: T10C1B1-38-TN

% Mức: NB
% ID: T10C1B1-40-TN
% ID: T10C1B1-40-TN
% Mức: NB
% ID: T10C1B1-39-TN
% ID: T10C1B1-39-TN
% Mức: NB
% ID: T10C1B1-39-TN

% Mức: NB
% ID: T10C1B1-41-TN
% ID: T10C1B1-41-TN
% Mức: NB
% ID: T10C1B1-40-TN
% ID: T10C1B1-40-TN
% Mức: NB
% ID: T10C1B1-40-TN

% Mức: NB
% ID: T10C1B1-42-TN
% ID: T10C1B1-42-TN
% Mức: NB
% ID: T10C1B1-41-TN
% ID: T10C1B1-41-TN
% Mức: NB
% ID: T10C1B1-41-TN

% Mức: NB
% ID: T10C1B1-43-TN
% ID: T10C1B1-43-TN
% Mức: NB
% ID: T10C1B1-42-TN
% ID: T10C1B1-42-TN
% Mức: NB
% ID: T10C1B1-42-TN

% Mức: TH
% ID: T10C1B1-44-TN
% ID: T10C1B1-44-TN
% Mức: TH
% ID: T10C1B1-43-TN
% ID: T10C1B1-43-TN
% Mức: TH
% ID: T10C1B1-43-TN

% Mức: NB
% ID: T10C1B1-45-TN
% ID: T10C1B1-45-TN
% Mức: NB
% ID: T10C1B1-44-TN
% ID: T10C1B1-44-TN
% Mức: NB
% ID: T10C1B1-44-TN

% Mức: NB
% ID: T10C1B1-46-TN
% ID: T10C1B1-46-TN
% Mức: NB
% ID: T10C1B1-45-TN
% ID: T10C1B1-45-TN
% Mức: NB
% ID: T10C1B1-45-TN

% Mức: NB
% ID: T10C1B1-47-TN
% ID: T10C1B1-47-TN
% Mức: NB
% ID: T10C1B1-46-TN
% ID: T10C1B1-46-TN
% Mức: NB
% ID: T10C1B1-46-TN

% Mức: NB
% ID: T10C1B1-48-TN
% ID: T10C1B1-48-TN
% Mức: NB
% ID: T10C1B1-47-TN
% ID: T10C1B1-47-TN
% Mức: NB
% ID: T10C1B1-47-TN

% Mức: NB
% ID: T10C1B1-49-TN
% ID: T10C1B1-49-TN
% Mức: NB
% ID: T10C1B1-48-TN
% ID: T10C1B1-48-TN
% Mức: NB
% ID: T10C1B1-48-TN

% Mức: NB
% ID: T10C1B1-50-TN
% ID: T10C1B1-50-TN
% Mức: NB
% ID: T10C1B1-49-TN
% ID: T10C1B1-49-TN
% Mức: NB
% ID: T10C1B1-49-TN

% Mức: TH
% ID: T10C1B1-51-TN
% ID: T10C1B1-51-TN
% Mức: TH
% ID: T10C1B1-50-TN
% ID: T10C1B1-50-TN
% Mức: TH
% ID: T10C1B1-50-TN

% ID: 0D1N1-1
% Mức: TH
% ID: T10C1B1-52-TN
% ID: T10C1B1-52-TN
% Mức: TH
% ID: T10C1B1-51-TN
% ID: T10C1B1-51-TN
% Mức: TH
% ID: T10C1B1-51-TN

% Mức: NB
% ID: T10C1B1-53-TN
% ID: T10C1B1-53-TN
% Mức: NB
% ID: T10C1B1-52-TN
% ID: T10C1B1-52-TN
% Mức: NB
% ID: T10C1B1-52-TN

% Mức: NB
% ID: T10C1B1-54-TN
% ID: T10C1B1-54-TN
% Mức: NB
% ID: T10C1B1-53-TN
% ID: T10C1B1-53-TN
% Mức: NB
% ID: T10C1B1-53-TN

% Mức: NB
% ID: T10C1B1-55-TN
% ID: T10C1B1-55-TN
% Mức: NB
% ID: T10C1B1-54-TN
% ID: T10C1B1-54-TN
% Mức: NB
% ID: T10C1B1-54-TN

% Mức: NB
% ID: T10C1B1-56-TN
% ID: T10C1B1-56-TN
% Mức: NB
% ID: T10C1B1-55-TN
% ID: T10C1B1-55-TN
% Mức: NB
% ID: T10C1B1-55-TN

% Mức: TH
% ID: T10C1B1-57-TN
% ID: T10C1B1-57-TN
% Mức: TH
% ID: T10C1B1-56-TN
% ID: T10C1B1-56-TN
% Mức: TH
% ID: T10C1B1-56-TN

% Mức: NB
% ID: T10C1B1-58-TN
% ID: T10C1B1-58-TN
% Mức: NB
% ID: T10C1B1-57-TN
% ID: T10C1B1-57-TN
% Mức: NB
% ID: T10C1B1-57-TN

% Mức: NB
% ID: T10C1B1-59-TN
% ID: T10C1B1-59-TN
% Mức: NB
% ID: T10C1B1-58-TN
% ID: T10C1B1-58-TN
% Mức: NB
% ID: T10C1B1-58-TN

% Mức: TH
% ID: T10C1B1-60-TN
% ID: T10C1B1-60-TN
% Mức: TH
% ID: T10C1B1-59-TN
% ID: T10C1B1-59-TN
% Mức: TH
% ID: T10C1B1-59-TN
% Mức: TH

%%%=============EX_4=============%%%
\dangbt{Mệnh đề với mọi, tồn tại}
\begin{ex}
	% ID: T10C1B1-59-TN
	% Mức: TH
	Mệnh đề ``$\exists x \in \mathbb{R}\colon x^2=9$'' khẳng định rằng
	\choice
		{Nếu $x$ là số thực thì $x^2=9$}
		{\True Có ít nhất một số thực mà bình phương của nó bằng $9$}
		{Bình phương của một số thực bằng $9$}
		{Chỉ có một số thực bình phương bằng $9$}
	\loigiai{
		Mệnh đề ``$\exists x \in \mathbb{R}\colon x^2=9$'' khẳng định rằng: Có ít nhất một số thực mà bình phương của nó bằng $9$.
		}
\end{ex}
% Mức: TH
% ID: T10C1B1-62-TN
% ID: T10C1B1-62-TN
% Mức: TH
% ID: T10C1B1-61-TN
% ID: T10C1B1-61-TN
% Mức: TH
% ID: T10C1B1-61-TN

% Mức: TH
% ID: T10C1B1-63-TN
% ID: T10C1B1-63-TN
% Mức: TH
% ID: T10C1B1-62-TN
% ID: T10C1B1-62-TN
% Mức: TH
% ID: T10C1B1-62-TN
% Mức: TH

%%%=============EX_6=============%%%
\begin{ex}
	% ID: T10C1B1-62-TN
	% Mức: TH
	Ký hiệu nào sau đây dùng để viết đúng mệnh đề ``$\sqrt{5}$ không phải là số hữu tỉ''?
	\choice
		{\True $\sqrt{5} \notin \mathbb{Q}$}
		{$\sqrt{5} \notin \mathbb{R}$}
		{$\sqrt{5} \neq \mathbb{R}$}
		{$\sqrt{5} \not \subset \mathbb{Q}$}
	\loigiai{
		Mệnh đề ``$\sqrt{5}$ không phải là số hữu tỉ'' được kí hiệu là ``$\sqrt{5} \notin \mathbb{Q}$''.
		}
\end{ex}
% Mức: TH
% ID: T10C1B1-64-TN
% ID: T10C1B1-64-TN
% Mức: TH
% ID: T10C1B1-63-TN
% ID: T10C1B1-63-TN
% Mức: TH
% ID: T10C1B1-63-TN
% Mức: TH

%%%=============EX_7=============%%%
\begin{ex}
	% ID: T10C1B1-63-TN
	% Mức: TH
	Cho mệnh đề $B\colon$ \lq\lq$\exists x \in \mathbb{R}, x^2 = 2$\rq\rq. Phát biểu bằng lời của mệnh đề $B$ là
	\choice
		{\True Có ít nhất một số thực mà bình phương của nó bằng $2$}
		{Bình phương của mỗi số thực đều bằng $2$}
		{Nếu $x$ là số thực thì $x^2$ bằng $2$}
		{Chỉ có một số thực mà bình phương của nó bằng $2$}
	\loigiai{
		Phát biểu bằng lời của mệnh đề $B\colon$ \lq\lq$\exists x \in \mathbb{R}, x^2 = 2$\rq\rq\, là \lq\lq Có ít nhất một số thực mà bình phương của nó bằng $2$\rq\rq.
		}
\end{ex}
\dangbt{Mệnh đề với mọi, tồn tại}
\begin{ex}
Khi phân tích các mệnh đề về số học, hãy nhận xét:
\choiceTF
{a) Với mọi số nguyên $k$, $k^2 - k$ chia hết cho 2.}
{b) Tồn tại số nguyên tố $p$ sao cho $p^2 - 1$ là số nguyên tố.}
{c) Với mọi số thực $x$, nếu $x > 1$ thì $x^2 > x$.}
{d) Tồn tại số tự nhiên $n$ sao cho $n^2 + 1$ chia hết cho 3.}
\loigiai{
a) Đúng. Ta có $k^2 - k = k(k-1)$. Đây là tích của hai số nguyên liên tiếp, nên luôn có một số chẵn. Do đó, tích $k(k-1)$ luôn chia hết cho 2.
b) Sai. Nếu $p$ là số nguyên tố, thì $p \ge 2$.
Nếu $p=2$, $p^2-1 = 2^2-1 = 3$, là số nguyên tố.
Nếu $p=3$, $p^2-1 = 3^2-1 = 8$, không là số nguyên tố.
Nếu $p > 2$, thì $p$ là số nguyên tố lẻ. Khi đó $p^2-1 = (p-1)(p+1)$. Vì $p$ là số lẻ, $p-1$ và $p+1$ là hai số chẵn liên tiếp. Do đó $p-1 \ge 2$ và $p+1 \ge 4$. Vì $p-1$ và $p+1$ đều là số chẵn, nên $(p-1)(p+1)$ là hợp số (có ít nhất hai thừa số chẵn). Vậy không tồn tại số nguyên tố $p$ sao cho $p^2-1$ là số nguyên tố.
c) Đúng. Nếu $x > 1$, ta nhân cả hai vế của bất đẳng thức $x > 1$ với $x$ (là số dương), ta được $x \cdot x > 1 \cdot x$, tức là $x^2 > x$.
d) Sai. Xét các trường hợp của $n$ khi chia cho 3:
Nếu $n \equiv 0 \pmod{3}$, thì $n^2+1 \equiv 0^2+1 \equiv 1 \pmod{3}$.
Nếu $n \equiv 1 \pmod{3}$, thì $n^2+1 \equiv 1^2+1 \equiv 2 \pmod{3}$.
Nếu $n \equiv 2 \pmod{3}$, thì $n^2+1 \equiv 2^2+1 \equiv 4+1 \equiv 5 \equiv 2 \pmod{3}$.
Trong mọi trường hợp, $n^2+1$ không chia hết cho 3.
}
\end{ex}

\dangbt{Mệnh đề với mọi, tồn tại}
\begin{ex}
Cho tập hợp $A = \{1, 2, 3, 4, 5\}$. Mệnh đề "$ \exists x \in A, x^2 - 5x + 6 = 0 $" đúng hay sai? Nếu đúng, hãy ghi giá trị của $x$ thỏa mãn (là số nguyên).
\shortans{2}
\loigiai{Ta cần tìm $x \in A$ sao cho $x^2 - 5x + 6 = 0$.
Giải phương trình bậc hai: $x^2 - 5x + 6 = 0 \Leftrightarrow (x-2)(x-3) = 0$.
Các nghiệm của phương trình là $x=2$ hoặc $x=3$.
Cả hai giá trị $x=2$ và $x=3$ đều thuộc tập hợp $A$.
Vì có ít nhất một giá trị $x \in A$ thỏa mãn điều kiện, mệnh đề là đúng.
Giá trị $x$ thỏa mãn là 2 (hoặc 3, nhưng chỉ cần ghi một giá trị).}
\end{ex}

\dangbt{Mệnh đề với mọi, tồn tại}
\begin{ex}
Cho mệnh đề $P(n)$: "$n^2 - 1$ chia hết cho 3".
a) Mệnh đề $\forall n \in \mathbb{N}^*, P(n)$ đúng hay sai? Giải thích.
b) Mệnh đề $\exists n \in \mathbb{N}^*, P(n)$ đúng hay sai? Giải thích.
\loigiai{
a) Mệnh đề $\forall n \in \mathbb{N}^*, P(n)$ là "Với mọi số tự nhiên $n$ khác 0, $n^2 - 1$ chia hết cho 3".
Mệnh đề này sai.
Giải thích: Chọn $n=2 \in \mathbb{N}^*$. Khi đó $n^2 - 1 = 2^2 - 1 = 4 - 1 = 3$. $3$ chia hết cho $3$.
Chọn $n=3 \in \mathbb{N}^*$. Khi đó $n^2 - 1 = 3^2 - 1 = 9 - 1 = 8$. $8$ không chia hết cho $3$.
Vì có một giá trị $n=3$ làm cho mệnh đề $P(n)$ sai, nên mệnh đề $\forall n \in \mathbb{N}^*, P(n)$ sai.

b) Mệnh đề $\exists n \in \mathbb{N}^*, P(n)$ là "Tồn tại số tự nhiên $n$ khác 0 sao cho $n^2 - 1$ chia hết cho 3".
Mệnh đề này đúng.
Giải thích: Chọn $n=1 \in \mathbb{N}^*$. Khi đó $n^2 - 1 = 1^2 - 1 = 0$. $0$ chia hết cho $3$.
Vì có ít nhất một giá trị $n=1$ làm cho mệnh đề $P(n)$ đúng, nên mệnh đề $\exists n \in \mathbb{N}^*, P(n)$ đúng.
}
\end{ex}
